\documentclass[11pt,a4paper]{article}
\usepackage[T1]{fontenc}
\usepackage[utf8]{inputenc}
\usepackage[spanish]{babel}
\usepackage{amsmath,amssymb}
\usepackage[margin=2.5cm]{geometry}
\usepackage[hidelinks]{hyperref}
\newcommand{\B}{\{0,1\}}
\title{Murdoku: modelo del caso de uso\\Generar una soluci\'on}
\author{Documentaci\'on de la implementaci\'on}
\date{}

\begin{document}
\maketitle

\section{Alcance y correspondencia con el c\'odigo}

El caso de uso \texttt{generate-solution}, ejecutado mediante
\texttt{murdoku-creator}, construye un mapa con habitaciones y muebles y
una colocaci\'on completa de personas. El resultado es una soluci\'on
objetivo $S^\star$, no un puzzle con pistas ni una prueba de unicidad.
El modelo optimiza primero el n\'umero de habitaciones amuebladas y,
como criterio secundario, la proximidad de las celdas a sus semillas.

La implementaci\'on se encuentra en:
\begin{itemize}
  \item \path{murdoku_checker\domain\generation_config.py}.
  \item \path{murdoku_checker\apps\solution_creator\solution_creator.py}:
        prepara la especificaci\'on y llama al modelo.
  \item \path{murdoku_checker\apps\solution_creator\placement_model.py}:
        construye y resuelve \texttt{PlacementModel}.
\end{itemize}

\section{Datos y preparaci\'on fuera de CP-SAT}

Sean $H,W\geq 1$ las dimensiones de la cuadr\'icula, $K\geq 1$ el
n\'umero de habitaciones, $T\geq 0$ el n\'umero de tipos de mueble y
$\alpha\in[0,100]$ el porcentaje de cobertura. Se definen
\[
 N=HW,\qquad n=\min(H,W),\qquad
 E=\operatorname{round}\left(N\frac{\alpha}{100}\right).
\]
El redondeo es el de Python: al entero m\'as cercano y, en un empate,
al entero par. $E$ cuenta celdas cubiertas, no ejemplares.
La configuraci\'on exige $K\leq N$, $T\leq E$ y, si $E>0$, $T>0$.
Estas validaciones no garantizan la factibilidad del modelo completo.

Antes de construir CP-SAT, \texttt{generate\_specification} fija para
cada tipo $t$:
\[
 s_t\in\{1,2\},\qquad m_t\in\mathbb Z_{\geq 1},\qquad o_t\in\B,
 \qquad \sum_{t=1}^{T}s_tm_t=E.
\]
$s_t$ es el tama\~no, $m_t$ el n\'umero de ejemplares y $o_t=1$
indica que las personas pueden ocupar sus celdas. Estos valores
\textbf{son datos del modelo}, no variables de decisi\'on.
Si $T=0$, la especificaci\'on de muebles es vac\'ia.

Se elige aleatoriamente qu\'e tipos tienen tama\~no dos, manteniendo
el presupuesto y dejando un tipo de tama\~no uno si $E$ es impar.
Cada tipo empieza con un ejemplar y se a\~naden ejemplares compatibles
con el presupuesto restante. Por defecto, la ocupabilidad tambi\'en
se sortea. La semilla controla estas decisiones.

La configuraci\'on admite \texttt{person\_names}, con exactamente
$n$ nombres, y \texttt{element\_names}, con exactamente $T$ entradas.
Cada lista debe contener nombres \'unicos y no vac\'ios, excluyendo
los formados solo por espacios. Si se omiten, se conservan
\texttt{person\_1}, etc., y \texttt{element\_1}, etc.
La \'ultima persona sigue siendo la v\'ictima.
Cada entrada de \texttt{element\_names} puede ser un nombre o un objeto
con \texttt{name} e \texttt{is\_occupiable}. Un objeto fija $o_t$
con un booleano JSON; si omite \texttt{is\_occupiable}, vale falso.
Una entrada de texto conserva el sorteo de ocupabilidad. Ambos
formatos pueden mezclarse. Los nombres y ocupabilidades se fijan
antes de resolver, no se cambian para recuperar factibilidad.

\section{Conjuntos e \'indices}

\begin{description}
  \item[$V=\{0,\ldots,N-1\}$] Celdas. Una celda de coordenadas $(i,j)$
        tiene \'indice $c=iW+j$, con $0\leq i<H$, $0\leq j<W$.
  \item[$\mathcal N(c)$] Vecinos ortogonales existentes de $c$:
        distancia Manhattan uno, sin diagonales.
  \item[$\mathcal K=\{0,\ldots,K-1\}$] Identificadores de habitaci\'on.
  \item[$s_k\in V$] Semilla fija de la habitaci\'on $k$. Se sortea la
        primera y se eligen las siguientes maximizando la distancia
        Manhattan a la semilla m\'as cercana, con desempates aleatorios.
  \item[$P=\{1,\ldots,n\}$] Personas. La v\'ictima es $v=n$.
  \item[$\mathcal I$] Ejemplares individuales de mueble. Cada ejemplar
        $i$ pertenece a un tipo $t(i)$.
  \item[$\mathcal A_i$] Huellas permitidas del ejemplar $i$. Si
        $s_{t(i)}=1$, todas las celdas individuales; si vale dos,
        todos los pares ortogonales no ordenados.
\end{description}

\section{Variables de decisi\'on}

\begin{align*}
 r_c&\in\mathcal K
     &&\text{habitaci\'on de la celda }c,\\
 d_c&\in\{0,\ldots,N-1\}
     &&\text{profundidad para demostrar conectividad},\\
 a_{cu}&\in\B \quad(u\in\mathcal N(c))
     &&\text{$c$ elige a $u$ como padre},\\
 f_{iA}&\in\B \quad(A\in\mathcal A_i)
     &&\text{el ejemplar $i$ ocupa la huella $A$},\\
 b_c&\in\B
     &&\text{la celda $c$ est\'a bloqueada},\\
 y_p&\in\{0,\ldots,H-1\},\quad
 x_p\in\{0,\ldots,W-1\}
     &&\text{fila y columna de la persona }p,\\
 q_p&\in V
     &&\text{\'indice lineal de su celda},\\
 \rho_p&\in\mathcal K
     &&\text{habitaci\'on de la persona }p,\\
 h_p&\in\B \quad(p\neq v)
     &&\text{$p$ comparte habitaci\'on con la v\'ictima}.
\end{align*}

\section{Restricciones}

\subsection{Partici\'on en habitaciones conexas}

Cada habitaci\'on tiene una \'unica ra\'iz fijada antes de resolver:
\begin{align}
 r_{s_k}&=k,\quad d_{s_k}=0 &&\forall k\in\mathcal K.
\end{align}
Sea $S=\{s_k:k\in\mathcal K\}$. Cada celda no ra\'iz elige
exactamente un vecino como padre:
\begin{align}
 \sum_{u\in\mathcal N(c)}a_{cu}&=1 &&\forall c\in V\setminus S,\\
 a_{cu}=1&\Longrightarrow r_c=r_u,\ d_c>d_u
     &&\forall c\in V\setminus S,\ u\in\mathcal N(c).
\end{align}

La profundidad decrece estrictamente al seguir padres, por lo que no
puede haber ciclos. Toda cadena termina en una ra\'iz de la misma
habitaci\'on. Como existe exactamente una ra\'iz por identificador,
cada habitaci\'on es no vac\'ia y conexa. No se impone un tama\~no
m\'inimo uniforme de cuatro celdas. La compacidad y la
proporci\'on se imponen mediante las restricciones siguientes.

\subsection{Como m\'aximo una habitaci\'on de una fila o columna}

Se introducen booleanos de pertenencia $m_{kc}$ y presencia por
coordenada $u_{ki}$, $w_{kj}$, todos con equivalencia exacta:
\begin{align}
 m_{kc}=1&\Longleftrightarrow r_c=k,\\
 u_{ki}&=\max_{\substack{c\in V\\\operatorname{fila}(c)=i}}m_{kc},\\
 w_{kj}&=\max_{\substack{c\in V\\\operatorname{columna}(c)=j}}m_{kc}.
\end{align}
As\'i, $\sum_i u_{ki}$ cuenta filas distintas ocupadas por $k$ y
$\sum_j w_{kj}$ cuenta columnas distintas. Las habitaciones ya son
no vac\'ias. Se definen tres booleanos por habitaci\'on:
\begin{align}
 s_k^{\mathrm{fila}}=1
   &\Longleftrightarrow \sum_i u_{ki}=1,\\
 s_k^{\mathrm{columna}}=1
   &\Longleftrightarrow \sum_j w_{kj}=1,\\
 t_k&=\max\{s_k^{\mathrm{fila}},s_k^{\mathrm{columna}}\}.
\end{align}
Finalmente se impone el l\'imite global
\[
 \boxed{\sum_{k\in\mathcal K}t_k\leq1.}
\]
Las dos orientaciones comparten el mismo l\'imite, no hay una
excepci\'on por eje. Una habitaci\'on de una sola celda activa ambos
$s_k$, pero cuenta una sola vez en $t_k$. Se permite que no exista
ninguna habitaci\'on estrecha. Todas las dem\'as ocupan al menos
dos filas y dos columnas.

\subsection{Formas compactas y proporci\'on m\'axima de 1.5}

Se introducen dimensiones enteras y un \'area por habitaci\'on:
\begin{align}
 H_k&=\sum_i u_{ki},&1\leq H_k&\leq H,\\
 W_k&=\sum_j w_{kj},&1\leq W_k&\leq W,\\
 A_k&=\sum_{c\in V}m_{kc},&1\leq A_k&\leq N.
\end{align}
La conectividad ya impuesta garantiza que las filas y columnas
ocupadas son consecutivas. Por ello, $H_k$ y $W_k$ son las
dimensiones del rect\'angulo envolvente. Se exige
\[
 \boxed{4A_k\geq3H_kW_k,\qquad
        2H_k\leq3W_k,\qquad 2W_k\leq3H_k.}
\]
El producto auxiliar se implementa con
\texttt{add\_multiplication\_equality}. La habitaci\'on ocupa al
menos el 75\% de su envolvente, permitiendo formas en L.
Adem\'as, para cada fila y columna, la secuencia de pertenencia
$m_{kc}$ debe pertenecer al lenguaje $0^*1^*0^*$, mediante
\texttt{add\_automaton}. As\'i se evitan huecos y formas en U.
Las dos desigualdades equivalen exactamente a
$\max(H_k,W_k)/\min(H_k,W_k)\leq3/2$, sin usar coma flotante.

Se permiten, por ejemplo, $2\times3$ y $3\times4$, pero no
$1\times2$ ni $2\times4$. El l\'imite de habitaciones estrechas
se conserva: solo permite como m\'aximo una habitaci\'on $1\times1$,
y puede no haber ninguna.
El auxiliar \texttt{\_build\_room\_shape\_constraints} construye
estas restricciones. Pueden hacer infactibles configuraciones
posibles con otras semillas.
No se cambia autom\'aticamente el mapa solicitado para evitarlas.

\subsection{Preferencia espacial aleatoria}

Con $\delta(c,s_k)$ la distancia Manhattan y
$\epsilon_{kc}\in\{0,\ldots,9\}$ un coste aleatorio fijo, se define
\[
 C=\sum_{k,c}(10\delta(c,s_k)+\epsilon_{kc})m_{kc}.
\]
La semilla de configuraci\'on controla tambi\'en estos costes.
Se favorecen habitaciones cercanas a sus semillas sin imponer
\'areas iguales. $C$ es el criterio secundario del objetivo descrito
m\'as abajo. La b\'usqueda se limita a diez segundos y puede
devolver una soluci\'on factible sin demostrar optimalidad.
Un resultado inconcluso sin soluci\'on se comunica como error.

\subsection{Colocaci\'on de los muebles}

Cada ejemplar elige una \'unica huella y dos muebles no se solapan:
\begin{align}
 \sum_{A\in\mathcal A_i}f_{iA}&=1 &&\forall i\in\mathcal I,\\
 \sum_{\substack{i\in\mathcal I,\ A\in\mathcal A_i\\c\in A}}f_{iA}
     &\leq 1 &&\forall c\in V.
\end{align}
Un mueble de dos celdas no puede atravesar una pared:
\[
 f_{i,\{c,u\}}=1\Longrightarrow r_c=r_u.
\]
El bloqueo se calcula mediante
\[
 b_c=
 \sum_{\substack{i\in\mathcal I,\ A\in\mathcal A_i\\
                   c\in A,\ o_{t(i)}=0}}f_{iA}
 \qquad\forall c\in V.
\]
La restricci\'on de no solapamiento hace que la suma sea booleana.
Las celdas sin muebles no est\'an bloqueadas. El presupuesto $E$ se
cumple por los tama\~nos y cantidades fijados, sin necesidad de una
restricci\'on adicional sobre el porcentaje.

\subsection{Preferencia por habitaciones amuebladas}

Se cuenta cualquier mueble, ocupable o bloqueado; las personas no
cuentan como objetos. Se introducen indicadores exactos
\begin{align}
 z_c&=\sum_{\substack{i\in\mathcal I,\ A\in\mathcal A_i\\c\in A}}f_{iA},
       &z_c&\in\B,\\
 e_{kc}&=\min(m_{kc},z_c),
       &e_{kc}&\in\B,\\
 g_k&=\max_{c\in V}e_{kc},
       &g_k&\in\B.
\end{align}
As\'i, $g_k=1$ si y solo si la habitaci\'on $k$ contiene al menos
una celda de mueble. Un mueble de dos celdas cuenta para una sola
habitaci\'on, porque no puede atravesar paredes.
El objetivo completo es
\[
 \boxed{\min\ C-M\sum_{k\in\mathcal K}g_k,\qquad
 M=N\bigl(10(H+W-2)+9\bigr)+1.}
\]
Cada celda pertenece exactamente a una habitaci\'on, por lo que
$0\leq C\leq M-1$. Una habitaci\'on amueblada adicional mejora el
objetivo m\'as que cualquier variaci\'on posible de $C$.
Esto equivale a maximizar lexicogr\'aficamente el n\'umero de
habitaciones amuebladas y despu\'es minimizar $C$.

Es una preferencia, no se exige $g_k=1$ para todo $k$. Se permiten
habitaciones sin muebles cuando faltan ejemplares o las piezas no
caben bajo las dem\'as restricciones. No se modifican cantidades,
tama\~nos ni ocupabilidad. Si no hay muebles, todos los indicadores
son cero y solo se minimiza $C$; la implementaci\'on omite entonces
los auxiliares de amueblamiento.

\subsection{Colocaci\'on de las personas}

\begin{align}
 q_p&=Wy_p+x_p &&\forall p\in P,\\
 b_{q_p}&=0 &&\forall p\in P,\\
 \operatorname{AllDifferent}\bigl((y_p)_{p\in P}\bigr),&\\
 \operatorname{AllDifferent}\bigl((x_p)_{p\in P}\bigr).&
\end{align}
La lectura $b_{q_p}=0$ es una restricci\'on \texttt{AddElement}.
No se repiten filas ni columnas entre personas. En particular, tampoco
pueden compartir celda. Una persona s\'i puede coincidir con un mueble
ocupable.

\subsection{Habitaci\'on de la v\'ictima}

Se obtiene la habitaci\'on de cada persona mediante otra lectura
indexada y se exige exactamente un acompa\~nante de la v\'ictima:
\begin{align}
 \rho_p&=r_{q_p} &&\forall p\in P,\\
 h_p=1&\Longleftrightarrow \rho_p=\rho_v &&\forall p\neq v,\\
 \sum_{p\in P\setminus\{v\}}h_p&=1.
\end{align}
Por tanto, la habitaci\'on de la v\'ictima contiene exactamente dos
personas. No se identifica una variable independiente de asesino.
Si $n=1$, la suma es vac\'ia y el modelo es infactible, aunque la
configuraci\'on haya superado las validaciones iniciales.

\section{Resoluci\'on y extracci\'on}

CP-SAT usa ocho trabajadores, \texttt{PORTFOLIO\_SEARCH},
\texttt{randomize\_search=True} y una semilla de solver derivada de la
semilla de configuraci\'on. La aleatorizaci\'on de la b\'usqueda no
equivale a muestrear uniformemente todas las soluciones factibles.

\begin{description}
  \item[\texttt{FEASIBLE} o \texttt{OPTIMAL}:] Se extrae una colocaci\'on
        factible, el mapa y cada ejemplar de mueble.
        \texttt{OPTIMAL} demuestra el \'optimo del objetivo espacial,
        no una mejor dificultad del puzzle.
        \texttt{FEASIBLE} no garantiza el m\'inimo de habitaciones
        sin muebles.
  \item[\texttt{INFEASIBLE}:] La app devuelve \texttt{None}; el comando
        informa del fallo. No cambia cantidades, tama\~nos ni ocupabilidad
        ni vuelve a sortear la especificaci\'on.
  \item[Otro estado:] Se produce \texttt{RuntimeError}; no se presenta
        como una demostraci\'on de imposibilidad.
\end{description}

El resultado contiene \texttt{GeneratedSolution}: especificaci\'on,
mapa, posiciones, v\'ictima y muebles individuales. Se conservan
los nombres configurados de personas y elementos; las habitaciones
se serializan como \texttt{room\_1}, etc.
Las coordenadas internas y las descripciones de pistas empiezan en cero.
El registro conserva la configuraci\'on efectiva y la semilla para
repetir las entradas de la ejecuci\'on. La b\'usqueda paralela y el
l\'imite de tiempo no garantizan la misma colocaci\'on en cada intento.

Se guarda \texttt{solution.html}, con posiciones visibles y una
pista informativa que identifica a la v\'ictima y recuerda que
estaba a solas con otra persona. Esa frase no introduce otra
restricci\'on ni participa en una selecci\'on de pistas.

\section{L\'imites de la garant\'ia}

Este caso de uso no selecciona pistas, no minimiza dificultad ni
demuestra unicidad. Las habitaciones, los muebles y las personas se
deciden conjuntamente. Sus variables auxiliares de conectividad no
forman parte del puzzle que se entrega al jugador. La verificaci\'on
posterior usa un mapa fijo y un modelo de colocaci\'on distinto,
descrito en el documento \texttt{verify-puzzle.tex}.

\end{document}
